% source: manti.01.ma, 1ère SM, devoir 1er semestre (Drive 0B6pfSJcEVng0QmczdXhhdnd6Mms). Ex.1 barycentre ; Ex.3 = WTQ5 Ex.7 (non repris)
\ex{Suites}{Qmcz}{Ex. 2}{9}
Soit $(U_n)_{n\in\mathbb{N}}$ la suite définie par $U_0=1$ et $U_{n+1}=3U_n-n+\dfrac12$.
\begin{enumerate}
\item Montrer que $(\forall n\in\mathbb{N})\ U_n\geq\dfrac n2$.
\item Étudier la monotonie de la suite $(U_n)_{n\in\mathbb{N}}$.
\item On pose $V_n=U_n-\dfrac12n$ pour tout $n\in\mathbb{N}$.
\begin{enumerate}
\item Montrer que $(V_n)_{n\in\mathbb{N}}$ est une suite géométrique en précisant sa raison $q$, et calculer $U_n$ en fonction de $n$.
\item Montrer que $U_0+U_1+\dots+U_{n-1}=\dfrac{3^n-1}{2}+\dfrac{n(n-1)}{4}$.
\end{enumerate}
\end{enumerate}
\fin
\ex{Applications}{Qmcz}{Ex. 4}{11}
Soit $f$ l'application de $\mathbb{R}^+\setminus\{4\}$ vers $\mathbb{R}^+$ définie par $f(x)=\dfrac{1}{\left(\sqrt x-2\right)^2}$.
\begin{enumerate}
\item Résoudre dans $\mathbb{R}^+\setminus\{4\}$ l'équation $f(x)=4$. Que peut-on en déduire ?
\item Soit $g$ la restriction de $f$ à l'intervalle $I=[0,4[$.
\begin{enumerate}
\item Montrer que $g$ est injective sur $I$.
\item Montrer que $g(I)\subseteq\left[\dfrac14,+\infty\right[$.
\item Montrer que $g$ est une bijection de $I$ vers $\left[\dfrac14,+\infty\right[$ et déterminer sa bijection réciproque $g^{-1}$.
\end{enumerate}
\end{enumerate}
\fin
